\chapter{Bacterial Culture}

\section*{What Is This Test?}
Bacterial culture is a foundational microbiologic technique in which a clinical specimen (blood, urine, sputum, CSF, wound exudate, tissue, or other body fluid) is inoculated onto nutrient-rich agar media and incubated under controlled conditions to promote the growth and identification of pathogenic bacteria. The culture proceeds through a sequence of steps: specimen collection and transport, primary inoculation onto appropriate selective and nonselective agar media, incubation at optimal temperature (typically 35--37\textdegree{}C) and atmosphere (aerobic, 5--10\% CO$_2$, microaerophilic, or anaerobic), macroscopic examination of colony morphology (size, shape, color, hemolysis pattern, odor), Gram stain of colonies, biochemical and enzymatic identification (catalase, coagulase, oxidase, indole, carbohydrate utilization), and identification by MALDI-TOF mass spectrometry, automated biochemical systems (VITEK 2, BD Phoenix, MicroScan), or nucleic acid-based methods. Definitive identification is followed by antimicrobial susceptibility testing (AST) to determine the minimum inhibitory concentration (MIC) of relevant antibiotics, using broth microdilution, disk diffusion (Kirby-Bauer), gradient diffusion (Etest), or automated systems. The entire process from specimen collection to final identification and AST typically requires 48--72 hours. Bacterial culture remains the gold standard for diagnosis of most bacterial infections, providing both identification of the causative pathogen and antimicrobial susceptibility to guide targeted therapy.

\section*{Alternative Names}
Bacteriology Culture; Culture and Sensitivity (C\&S); Aerobic Culture; Anaerobic Culture; Microbiologic Culture; Routine Culture

\section*{Principle}
The core principle of bacterial culture is the provision of an artificial environment containing all the nutritional, osmotic, pH, and atmospheric conditions required for bacterial replication. Each specimen type is plated onto a panel of media selected for that specific specimen and the organisms most likely to be pathogens at that site. Blood agar (tryptic soy agar with 5\% sheep blood): General purpose growth medium for most bacteria and fungi; allows detection of hemolysis patterns (alpha = partial/green, beta = complete/clear, gamma = none). Chocolate agar: Blood agar heated to 80\textdegree{}C to lyse RBCs and release hemin (X factor) and NAD (V factor), supporting fastidious organisms (\textit{Haemophilus influenzae}, \textit{Neisseria gonorrhoeae}, \textit{N. meningitidis}). MacConkey agar: Selective (bile salts, crystal violet inhibit gram-positive) and differential (lactose fermentation turns colonies pink---lactose fermenters like \textit{E. coli}, \textit{Klebsiella}; non-fermenters like \textit{Pseudomonas}, \textit{Acinetobacter} remain colorless). Colistin-nalidixic acid (CNA) or phenylethyl alcohol (PEA) agar: Selective for gram-positive organisms. Anaerobic media (CDC anaerobe blood agar, Kanamycin-vancomycin laked blood agar, phenylethyl alcohol agar): Prereduced media incubated in anaerobic jars or chambers (N$_2$/H$_2$/CO$_2$ atmosphere with palladium catalyst). Selective broths (thioglycollate broth, Lim enrichment broth for GBS): Liquid media to enhance recovery of small numbers of organisms or fastidious species. Incubation: 35--37\textdegree{}C for 18--48 hours (extended to 7--14 days for certain organisms). Colony identification by MALDI-TOF mass spectrometry involves smearing a colony on a target plate, overlaying with matrix (alpha-cyano-4-hydroxycinnamic acid), and analyzing the protein mass spectrum against a reference database. MALDI-TOF provides species-level identification in minutes with accuracy comparable to 16S rRNA gene sequencing at a fraction of the cost and time.

\section*{History}
Robert Koch pioneered solid culture media in the 1880s using potato slices and gelatin, later improved by his assistant Walther Hesse's wife, Fanny Hesse, who suggested agar (a seaweed-derived polysaccharide she used for fruit preserves) as a solidifying agent. The Petri dish was invented by Julius Petri in 1887. Selective and differential media were developed in the late 19th and early 20th centuries. The Kirby-Bauer disk diffusion method was standardized in 1966 \cite{bauer1966antibiotic}. Automated blood culture systems (BACTEC, BacT/ALERT) were introduced in the 1970s--1980s, detecting CO$_2$ production by growing organisms. Automated identification and susceptibility systems (VITEK, MicroScan) were introduced in the 1980s. MALDI-TOF mass spectrometry revolutionized bacterial identification in the 2010s, reducing the time to species identification from 24--48 hours to minutes. Despite advances in molecular diagnostics (multiplex PCR, metagenomic next-generation sequencing), culture remains irreplaceable for antimicrobial susceptibility testing, epidemiologic typing, and the recovery of organisms for whole-genome sequencing that informs outbreak investigation and molecular epidemiology.

\section*{Why This Test Is Ordered}
\begin{itemize}
  \item Diagnosis of the causative bacterial pathogen in patients with suspected infection: bloodstream infection (sepsis/bacteremia), pneumonia, urinary tract infection, wound/surgical site infection, meningitis, peritonitis, septic arthritis, osteomyelitis, endocarditis, and abscesses
  \item Isolation and identification of the specific bacterial species so that targeted (narrow-spectrum) antimicrobial therapy can be administered, replacing broad-spectrum empiric therapy
  \item Antimicrobial susceptibility testing (AST) of the isolated pathogen to guide definitive antibiotic selection and detect resistance mechanisms (MRSA, ESBL, carbapenemase, VRE, multidrug-resistant \textit{Pseudomonas/ Acinetobacter})
  \item Recovery and storage of the bacterial isolate for further testing: molecular typing (PFGE, MLST, WGS) for outbreak investigation, detection of specific resistance genes, and epidemiologic surveillance (NHSN, CDC)
  \item Screening cultures (active surveillance) for multidrug-resistant organisms: MRSA nasal swab, VRE rectal swab, carbapenem-resistant Enterobacterales (CRE) rectal swab, carbapenem-resistant \textit{Acinetobacter} (CRAB), and \textit{Candida auris} axilla/groin swab
  \item Evaluation of sterile body sites (CSF, synovial fluid, pleural fluid, peritoneal fluid) where any bacterial growth is potentially significant and requires identification and susceptibility testing
\end{itemize}

\section*{Is This a Primary or Secondary Test}
Bacterial culture is the primary test for the diagnosis of bacterial infections. It is the gold standard against which newer molecular diagnostic methods are measured. Culture is indispensable for antimicrobial susceptibility testing; no molecular test currently replaces phenotypic AST.

\section*{How This Test Is Performed}
Specimen collection is the most critical step. Detailed collection instructions for each specimen type are provided in the respective chapters (blood culture, sputum culture, urine culture, etc.). Key principles: (1) Collect specimen BEFORE antibiotic administration whenever possible. (2) Avoid contamination with commensal flora: use aseptic technique for normally sterile sites; instruct patients on proper midstream urine and deep-cough sputum collection. (3) Collect an adequate volume: blood cultures require 20 mL/set in adults; CSF requires at least 0.5--1 mL per tube. (4) Transport promptly: most specimens should reach the laboratory within 2 hours at room temperature. CSF, sterile body fluids, and tissue should be transported STAT and never refrigerated (unless specified for a particular organism). (5) Use appropriate transport media: swabs should be placed in Amies transport medium with charcoal; specimens for anaerobic culture require prereduced anaerobic transport containers. (6) Label specimens completely: patient identification, date/time of collection, anatomic site, and specific requests (e.g., "rule out Nocardia," "stat Gram stain").

\section*{What Preparation Is Needed}
Preparation is specimen-specific. Key principles: (a) surgical wound and sterile body site preparation with antiseptic (chlorhexidine, povidone-iodine) to minimize contamination; (b) mouth rinse with water before sputum collection; (c) cleansing of the urethral meatus before midstream urine collection; (d) withholding of antibiotics until cultures are obtained, when clinically safe. Fasting is not required for routine cultures.

\section*{Contraindications and Precautions}
The most common error in culture interpretation is attributing clinical significance to a contaminant. Blood culture contaminated with coagulase-negative staphylococci, \textit{Corynebacterium} species, \textit{Bacillus} species, or \textit{Cutibacterium} (Propionibacterium) species from skin flora is common (2--3\% contamination rate). Growth of these organisms in a single blood culture set most likely represents contamination; growth in multiple sets increases the likelihood of true infection. Sputum specimens contaminated with oropharyngeal flora (presence of >10 squamous epithelial cells per LPF) should be rejected or reported with a caveat. Urine cultures with mixed growth of $\geq$3 organisms suggest contamination. Anaerobic cultures require specific transport media and plating within 15--30 minutes of collection; delays result in loss of obligate anaerobes. Prior antibiotic therapy dramatically reduces culture yield. If empiric antibiotics must be administered before cultures (e.g., suspected meningococcemia, septic shock), the cultures should be obtained as soon as possible after antibiotics (cultures remain positive for hours after the first antibiotic dose in most cases, but sensitivity is reduced).

\section*{Risks}
Risks are associated with specimen collection (venipuncture, lumbar puncture, biopsy) rather than the culture process. The major clinical risk is a false-negative culture (sterile culture in a patient with a true infection) leading to delayed or inappropriate antibiotic therapy. Conversely, interpreting a contaminated culture as representing true infection leads to unnecessary and potentially harmful antibiotic exposure.

\section*{What Happens After the Test}
Preliminary Gram stain results are available within 1--2 hours of specimen receipt. Bacterial growth is typically detected within 18--24 hours (aerobic), 24--48 hours (anaerobic), or 5--7 days (fastidious organisms). Final identification and AST results are available at 48--72 hours. Positive blood cultures are flagged as critical results; Gram stain result is communicated to the provider immediately. Empiric broad-spectrum antibiotics should be narrowed (de-escalated) based on culture results to target the identified pathogen(s) and their susceptibility profile.

\section*{Understanding Results}

\subsection*{Below Normal}
\begin{itemize}
  \item \textbf{No growth at 48 hours (or at final report):} No bacteria isolated. Interpretation depends on the specimen type. From normally sterile sites (blood, CSF, synovial fluid), this strongly argues against bacterial infection. From nonsterile sites (sputum, wound swabs), "no growth" or "normal flora" indicates that either no pathogens are present, or the specimen was inadequate
\end{itemize}

\subsection*{Above Normal}
\begin{itemize}
  \item \textbf{One or more pathogens isolated:} Organism(s) identified to the species level. Quantitative or semi-quantitative growth reported (e.g., >100,000 CFU/mL for clean-catch urine; >10$^5$ CFU/mL for BAL; moderate or many for sputum). Each pathogen's antimicrobial susceptibility profile is reported: susceptible (S), susceptible dose-dependent (SDD), intermediate (I), or resistant (R). MICs (mcg/mL) are provided for most antibiotics
  \item \textbf{Mixed growth or polymicrobial:} Growth of $\geq$3 organisms suggests contamination (urine, sputum) unless the specimen is from a non-sterile polymicrobial site (abscess, peritonitis, diabetic foot wound), where mixed flora is expected and may include aerobes and anaerobes
\end{itemize}

\section*{Normal Results}
Sterile body sites (blood, CSF, synovial fluid, pleural fluid, peritoneal fluid, pericardial fluid): \textbf{No growth at 5 days.} Nonsterile sites (sputum, urine, wounds): \textbf{Normal flora} or \textbf{No significant growth} or \textbf{Negative for pathogens.}

\section*{Follow-up Tests}
\begin{itemize}
  \item \textbf{Antimicrobial susceptibility testing (AST):} Performed on all clinically significant isolates from normally sterile sites and selected nonsterile sites. Broth microdilution MIC is the gold standard \cite{jorgensen2009antimicrobial}. Results are interpreted using CLSI or EUCAST breakpoints
  \item \textbf{Blood cultures (repeat):} For patients with bacteremia, blood cultures should be repeated at 48--72 hours to document clearance (sterile after 48--96 hours is expected for uncomplicated bacteremia with appropriate therapy). Persistent bacteremia suggests an uncontrolled source, endocarditis, or antibiotic resistance
  \item \textbf{Identification of resistance mechanisms:} Molecular testing for resistance genes (\textit{mecA} for MRSA, \textit{vanA/vanB} for VRE, \textit{bla}$_{CTX-M}$/\textit{bla}$_{KPC}$/\textit{bla}$_{NDM}$ for beta-lactamase-mediated resistance) may be performed to augment or confirm phenotypic AST results
  \item \textbf{16S rRNA gene sequencing:} When an organism cannot be identified by standard methods (MALDI-TOF, biochemical). PCR amplification and sequencing of conserved 16S ribosomal RNA gene regions identifies the organism to the genus or species level
  \item \textbf{Bacterial isolate storage:} Bacterial isolates are stored frozen (-70\textdegree{}C) in glycerol or bead storage systems for future testing (outbreak investigation, retrospective AST, whole-genome sequencing)
\end{itemize}

\section*{References}
\bibliographystyle{unsrtnat}
\bibliography{references}

\clearpage
